% -*- mode : latex; coding: utf-8 -*-
\chapter{Optimization Design} \label{sec:Optimization}
The underlying source of inefficiency in current state-of-the-art RH mitigation schemes is that MC doesn't know
the mapping between physical to logical rows. Due to rapid technology scaling in DRAM foundry, current DRAM devices
may have faulty cells due to manufacturing imperfection. DRAM vendors correct such errors by replacing rows containing 
faulty cells with fault-free rows. This approach is often called row sparing and has been widely adopted by vendors. 
Each DRAM bank has spare rows, and when faulty cells are found at the test phase after the manufacturing process, test program \textit{remaps} 
faulty row's address to fault-free spare row and saves remapping information into programmable memory (remapping table). At run-time, 
the global row decoder consults the remapping table and finds the corresponding row. Also, the majority of DRAM devices use a simple hash 
function between logical row address (sent by MC) and physical row address (actual row layout) and this hash function 
differ between vendors and even products within the same vendor. 
All the above schemes complicate MC to find out actual adjacent rows of aggressor row, making the RH mitigation scheme inefficient. 
In this paper, I redesigned the state-of-the-art mitigation systems assuming MC knows the mapping between physical to logical rows.

\section{Graphene}
\subsection{Algorithm $\And$ Modification}
\label{sec:Graphene Modification}
As explained in \ref{sec:Existing Solutions} Graphene uses CAM counter (per bank) to track activation count for 
row adresses during the last \textbf{W} ACTs. When the counter reaches \textbf{T}, Graphene control logic signals mitigator to send NRR
for aggressor row. In order to implement the above scheme, the CAM counter must track $N_{entry}$ entries that follows inequality 
\ref{eqn:Graphene_entrynum} \cite{Graphene}:
\begin{equation}
\mathrm{N_{entry} >\frac{W}{T} - 1} \\
\label{eqn:Graphene_entrynum} 
\end{equation}
There are two things to consider when setting T at Graphene. First, Graphene refreshes its CAM counter (counter refresh) at every tREFW. 
This counter refresh is asynchronous to individual rows, so a maximum  $2(T-1)$ number of ACTs can go through Graphene without invoking 
mitigation logic (T-1 ACTs between row refresh and counter refresh, another T-1 ACTs between counter refresh and next row refresh). 
Also, since every ACTs is considered as part of a double-hammering attack pattern, Graphene must use a threshold of $Flip_{th}/2$. 
Considering both points in the above discussion, we get the following inequality \ref{eqn:Graphene_T}.
\begin{equation}
\mathrm{2(T-1) <\frac{Flip_{th}}{2} \implies T<\frac{Flip_{th}}{4}+1} \\
\label{eqn:Graphene_T} 
\end{equation}
If MC is extended to know logical to physical row mapping, we can track victim rows instead of aggressor rows. 
The optimized version of Graphene (\textit{Graphene-op}) uses the same algorithm, but instead of focusing on the aggressor, it determines
the victims of ACT command, updates counter, and sends \textit{refresh} command instead of NRR. We don't have to consider every ACTs 
as part of the double-hammering pattern, instead of using a halved threshold. However, since we always have to count two victims 
for a single ACT, the window size is doubled. So we get inequality \ref{eqn:Graphene-op} for Graphene-op.

\begin{equation}
\label{eqn:Graphene-op} 
\begin{gathered}
\mathrm{2(T-1) <Flip_{th} \implies T<\frac{Flip_{th}}{2}+1} \\
\mathrm{N_{entry} >\frac{2W}{T} - 1} \\
\end{gathered}
\end{equation}
Note that although the number of tracking elements is doubled, T is also doubled, making Graphene-op's CAM counter entry ($N_{entry}$) 
the same as Graphene. In addition, a Half-double attack is impossible in Graphene-op since it doesn't use NRR command. 
In conclusion, Graphene-op uses a threshold twice the size of Graphene without incurring almost no storage overhead and 
effectively blocks Half-double attacks. 

\subsection{Security Analysis of Modified Algorithm}
Here I show that Graphene-op successfully mitigates RH by following the same flow as Graphene. 
I start by proving the following theorem:

\noindent\textbf{Theorem.} \textit{Tracking victim rows of last W aggressive ACTs in size of $N_{entry}$ table and refreshing entry whenever 
it becomes multiple of T prevents any row's actual count from reaching T without incurring refresh} 
\footnote{actual count for a given row refers to the number of ACTs sent to any adjacent aggressor rows between refresh command}  

\noindent\textbf{Lemma 1.} The (estimated) count of every entry in the table is equal to or greater than the actual count of any victim row.

\noindent\textbf{Lemma 2.} Spillover count cannot exceed $\frac{2W}{N_{entry}+1}$

\noindent\textbf{Proof of Lemma 1.}
By algorithm introduced in section \ref{sec:Existing Solutions} and modification in section \ref{sec:Graphene Modification}, the only way
an estimated count can become smaller than an actual count is when it is evicted. When it is evicted, its estimated count 
must be equal to the spillover count. Between eviction and reentering, it will update the spillover counter. 
When it is written back, it will set its counter to spillover counter+1. $\hfill \square$

\noindent\textbf{Proof of Lemma 2.}
Since every arrived ACTs increments the sum of the spillover counter and counters in the table by 2, 
the sum of the spillover counter and counters in the table has a maximum of $2W$. Since the spillover counter should be smaller or equal
to a minimum value in the counters in the table, it will have a maximum value when all entries in the table have the same value as 
the spillover counter: $\frac{2W}{N_{entry}+1}$ $\hfill \square$

\noindent\textbf{Proof of Theorem.} 
When the actual count of victim row X turns T-1 to T, row X's estimated count should be T or higher due to Lemma 1. So, refresh action 
must be performed or already has been performed in the past. By inequality \ref{eqn:Graphene-op} and Lemma 2,
the spillover counter is always smaller than T and thus the entry is never evicted from that moment. 
So, when another T number of ACTs affects row X, the estimated counter would have increased by exactly T amount, 
and refresh action is performed accordingly. $\hfill \square$

\section{Hydra}
\subsection{Algorithm $\And$ Modification}
Recent studies \cite{DeeperLook} showed that due to manufacturing factors and design factors in DRAM foundries, 
$Flip_{th}$ of each row differs greatly. In an optimized version of Hydra (Hydra-op), I encoded the unique $Flip_{th}$ of each row 
assuming it can be known at the test phase of the DRAM manufacturing process and embedded those metadata along with the row's RCT counter. 
Since RCT size is less than 1\% of typical DRAM, this method incurs a small storage overhead. 
There is also little performance degradation, as each row's counter and metadata can be placed in the same cache line. 
Although tracking victim rows is orthogonal to adding threshold metadata for each row in the DRAM counter, I did not implement it in the 
current paper due to its obvious speedup. However, current work can easily be extended to incorporate both optimizations. 

